\documentclass[11pt]{article}
\usepackage[T1]{fontenc}
\usepackage{amsmath,amssymb,amsthm,lmodern}
\usepackage[margin=1in]{geometry}
\usepackage{hyperref}
\hypersetup{hidelinks,pdftitle={Milne AV II Theorem 1.1: Frobenius and point counts},pdfauthor={J.S. Milne}}
\DeclareMathOperator{\AVEnd}{End}
\DeclareMathOperator{\AVHom}{Hom}
\DeclareMathOperator{\AVTr}{Tr}
\DeclareMathOperator{\AVPic}{Pic}
\DeclareMathOperator{\AVTgt}{Tgt}
\DeclareMathOperator{\AVdiv}{div}
\DeclareMathOperator{\AVProj}{Proj}
\newcommand{\AVQ}{\mathbb{Q}}
\newcommand{\AVZ}{\mathbb{Z}}
\newcommand{\AVR}{\mathbb{R}}
\newcommand{\AVC}{\mathbb{C}}
\newcommand{\AVF}{\mathbb{F}}
\newcommand{\AVP}{\mathbb{P}}
\newcommand{\AVA}{\mathbb{A}}
\newcommand{\AVid}{\operatorname{id}}
\newcommand{\AVsrc}[1]{\par\smallskip\noindent\textit{Source: #1.}\par\smallskip}
\newcommand{\AVeditorial}[1]{\par\smallskip\noindent\textsf{Editorial note.} #1\par\smallskip}
\title{Frobenius and the Point Counts of an Abelian Variety\\Chapter II, Theorem 1.1}
\author{J.S. Milne}
\date{Private reading transcription of AV v2.00 (2008)\\with separate supporting material from the corrected Storrs article, January 2, 2022}
\begin{document}
\maketitle

% BEGIN CONTEXT AV:II1.1
\noindent\textbf{Fixed sources and selection.}
The primary source is J.S. Milne, \emph{Abelian Varieties}, course notes,
172 physical pages. Its cover reads Version 2.0, March 16, 2008;
its bibliographic version is 2.00 and its internal history also records
March 17, 2008. These are the printed labels of the fixed copy, not an
inferred replacement version. The source URL is
\url{https://www.jmilne.org/math/CourseNotes/AV.pdf}, with landing page
\url{https://www.jmilne.org/math/CourseNotes/av.html}.
The primary selection is Chapter II, Section 1, physical pages 81--84
(printed pages 75--78), through the complete proofs of Lemmas 1.2 and 1.3.
Remark 1.4 is retained as uncounted closing context. The zeta-function
corollary following it and the Honda--Tate section are outside this selection.

The separate supporting source is Milne's corrected author file for
\emph{Abelian Varieties}, originally Chapter V of \emph{Arithmetic Geometry}
(Storrs, Conn., 1984), pp. 103--150, Springer, New York, 1986.
The actual fixed corrected PDF is dated \textbf{January 2, 2022} and has
49 physical pages. Its URL is
\url{https://www.jmilne.org/math/xnotes/AVs.pdf}, with author page
\url{https://www.jmilne.org/math/xnotes/AVs.html}.
Only Section 17's necessary context, Theorem 17.3, Lemma 17.4 and their
complete printed argument on physical/printed pages 35--36 are appended.
The 1986 publication, this corrected 2022 file, and the 2008 course notes
remain separately identified sources.

\AVeditorial{There is one primary target, II Theorem 1.1, for arbitrary
abelian varieties over a finite field, including non-simple ones, and all
integers $m\geq1$. No restriction excluding characteristics 2 or 3 is added.
II Lemmas 1.2 and 1.3 are retained as substantive steps of the author's
primary proof. I Proposition 10.20, I Corollary 10.24 and Rosati positivity
are explicit strong supporting inputs. Supporting statements, definitions
and the separately sourced Rosati calculation are not additional primary
targets. No new proof is supplied.}

\subsection*{Necessary notation and characteristic-polynomial inputs from the 2008 notes}
\AVsrc{I, Section 10, physical page 49, printed page 43; notation excerpt}
We let $\AVEnd^0(A)=\AVEnd(A)\otimes\AVQ$.

\AVsrc{I, Section 10, physical page 50, printed page 44; opening paragraph}
Let $A$ be an abelian variety of dimension $g$ over a field $k$.
Recall that, for any $m$ not divisible by the characteristic of $k$,
$A_m(k^{\mathrm{sep}})$ has order $m^{2g}$, and that, for a prime
$\ell\ne\operatorname{char}(k)$,
$T_\ell A=\varprojlim A_{\ell^n}(k^{\mathrm{sep}})$ (see 7.3).

\AVsrc{I Proposition 10.5 and following paragraph, physical page 51, printed page 45}
\noindent\textbf{Proposition 10.5.}
\emph{For $\ell\ne\operatorname{char}(k)$, $T_\ell A$ is a free
$\AVZ_\ell$-module of rank $2g$.}

A homomorphism $\alpha:A\to B$ induces a homomorphism
$A_n(k^{\mathrm{sep}})\to B_n(k^{\mathrm{sep}})$, and hence a homomorphism
\[
T_\ell\alpha:T_\ell A\longrightarrow T_\ell B,\qquad
(a_1,a_2,\ldots)\longmapsto(\alpha(a_1),\alpha(a_2),\ldots).
\]
Therefore $T_\ell$ is a functor from abelian varieties to $\AVZ_\ell$-modules.

\AVsrc{I, Section 10, physical page 52, printed page 46; rational Tate-space definition}
Write $V_\ell A=T_\ell A\otimes_{\AVZ_\ell}\AVQ_\ell
\;(=T_\ell A\otimes_{\AVZ}\AVQ)$.

\AVsrc{I Theorem 10.9 and its immediately preceding degree convention,
physical page 52, printed page 46}
Let $\alpha\in\AVEnd(A)$. If $\alpha$ is an isogeny, we define
$\deg(\alpha)$ as before; otherwise, we set $\deg(\alpha)=0$.

\noindent\textbf{Theorem 10.9.}
\emph{Let $\alpha\in\AVEnd(A)$. There is a unique monic polynomial
$P_\alpha\in\AVZ[X]$ of degree $2g$ such that
$P_\alpha(r)=\deg(\alpha-r)$ for all integers $r$.}

\AVsrc{I Remark 10.11, physical page 53, printed page 47; degree and polynomial extension}
For $\alpha\in\AVEnd(A)$ and $n\in\AVZ$,
\[
\deg(n\alpha)=\deg(n_A)\cdot\deg(\alpha)=n^{2g}\deg(\alpha).
\]
We can use this formula to extend the definition of $\deg$ to
$\AVEnd^0(A)$. Since $\AVEnd(A)$ is torsion-free, we can identify
$\AVEnd(A)$ with a submodule of $\AVEnd^0(A)$. For
$\alpha\in\AVEnd^0(A)$, define
\[
\deg(\alpha)=n^{-2g}\deg(n\alpha)
\]
if $n$ is any integer such that $n\alpha\in\AVEnd(A)$.
The previous formula shows that this is independent of the choice of $n$.
Similarly, once we have proved the theorem, we can define
\[
P_\alpha(X)=n^{-2g}P_{n\alpha}(nX),\qquad
\alpha\in\AVEnd^0(A),\quad n\alpha\in\AVEnd(A).
\]
Then $P_\alpha(X)$ is a monic polynomial of degree $2g$ with rational
coefficients, and
\[
P_\alpha(r)=\deg(\alpha-r),\qquad\text{any }r\in\AVQ.
\]

\AVsrc{I, Section 10, physical page 54, printed page 48; trace convention}
We call $P_\alpha$ the \emph{characteristic polynomial} of $\alpha$ and
we define the \emph{trace} $\AVTr(\alpha)$ of $\alpha$ by the equation
\[
P_\alpha(X)=X^{2g}-\AVTr(\alpha)X^{2g-1}+\cdots+\deg(\alpha).
\]

% BEGIN SUPPORT STATEMENT AV:I10.20
\AVsrc{I Proposition 10.20, physical page 56, printed page 50}
\noindent\textbf{Proposition 10.20.}
\emph{For all $\ell\ne\operatorname{char}(k)$, $P_\alpha(X)$ is the
characteristic polynomial of $\alpha$ acting on $V_\ell A$; hence the trace
and degree of $\alpha$ are the trace and determinant of $\alpha$ acting
on $V_\ell A$.}
% END SUPPORT STATEMENT AV:I10.20

% BEGIN SUPPORT STATEMENT AV:I10.24
\AVsrc{I Corollary 10.24, physical page 59, printed page 53}
\noindent\textbf{Corollary 10.24.}
\emph{Let $\alpha\in\AVEnd^0(A)$, and assume $\AVQ[\alpha]$ is a
product of fields. Let $C_\alpha(X)$ be the characteristic polynomial of
$\alpha$ acting on $\AVQ[\alpha]$ (e.g., if $\AVQ[\alpha]$ is a field,
this is the minimum polynomial of $\alpha$); then}
\[
\{a\in\AVC\mid C_\alpha(a)=0\}
=\{a\in\AVC\mid P_\alpha(a)=0\}.
\]
% END SUPPORT STATEMENT AV:I10.24

\subsection*{Polarization and the parent Rosati input}
\AVsrc{I, Section 11, physical page 59, printed page 53; polarization definition}
A \emph{polarization}\footnote[15]{This notion of polarization differs
slightly from Weil's original definition.}
$\lambda$ of an abelian variety $A$ is an isogeny
$\lambda:A\to A^{\vee}$ such that, over $k^{\mathrm{al}}$, $\lambda$
becomes of the form $\lambda_{\mathcal L}$ for some ample sheaf
$\mathcal L$ on $A_{k^{\mathrm{al}}}$. Unfortunately, this is not quite
the same as requiring that $\lambda$ itself be of the form
$\lambda_{\mathcal L}$ for $\mathcal L$ an ample invertible sheaf on $A$
(Milne 1986, 13.2).

\AVsrc{I, Section 14, physical page 67, printed page 61; Rosati definition}
Fix a polarization $\lambda$ on $A$. As $\lambda$ is an isogeny
$A\to A^{\vee}$, it has an inverse in
$\AVHom^0(A^{\vee},A)\overset{\mathrm{df}}=
\AVHom(A^{\vee},A)\otimes\AVQ$.
The \emph{Rosati involution} on $\AVEnd(A)\otimes\AVQ$ corresponding
to $\lambda$ is
\[
\alpha\longmapsto\alpha^{\dagger}
=\lambda^{-1}\circ\alpha^{\vee}\circ\lambda.
\]

% BEGIN SUPPORT PARENT ROSATI AV:I14.3
\AVsrc{I Theorem 14.3 and its complete printed proof, physical page 68, printed page 62}
The next result is very important.

\noindent\textbf{Theorem 14.3.}
\emph{The bilinear form}
\[
(\alpha,\beta)\longmapsto\AVTr(\alpha\circ\beta^{\dagger}):
\AVEnd(A)\otimes\AVQ\times\AVEnd(A)\otimes\AVQ\longrightarrow\AVQ
\]
\emph{is positive definite, i.e., $\AVTr(\alpha\circ\alpha^{\dagger})>0$
for $\alpha\ne0$. More precisely, let $D$ be the ample divisor defining
the polarization used in the definition of $\dagger$; then}
\[
\AVTr(\alpha\circ\alpha^{\dagger})
=\frac{2g}{(D^g)}(D^{g-1}\cdot\alpha^*(D)).
\]
\begin{proof}
As $D$ is ample and $\alpha^*(D)$ is effective, the intersection number
$(D^{g-1}\cdot\alpha^*(D))$ is positive. Thus the second statement implies
the first. The second statement is proved by a calculation, which we omit
for the present (see Milne 1986, \S17).
\end{proof}
% END SUPPORT PARENT ROSATI AV:I14.3

\AVeditorial{The preceding omission is part of the 2008 source and is
retained literally. The separate 2022 author calculation appended below
supports precisely this trace/intersection identity; it is not represented
as a calculation printed in the 2008 parent. The unrelated I13.12
polarization issue is not part of this selection. The statements of
I10.20 and I10.24 are retained as explicit inputs; their full preceding
chapter is not transcribed.}

\subsection*{Chapter II. Abelian Varieties: Arithmetic}
\subsubsection*{1. The Zeta Function of an Abelian Variety}
\AVsrc{2008 notes, physical page 81, printed page 75; complete introduction to II1.1}
We write $\AVF_q$ for a finite field with $q$ elements, $\AVF$ for an
algebraic closure of $\AVF_q$, and $\AVF_{q^m}$ for the unique subfield
of $\AVF$ with $q^m$ elements. Thus the elements of $\AVF_{q^m}$ are the
solutions of $c^{q^m}=c$.

For a variety $V$ over $\AVF_q$, the Frobenius map $\pi_V:V\to V$ is
defined to be the identity map on the underlying topological space of $V$
and is the map $f\mapsto f^q$ on $\mathcal O_V$. For example, if
$V=\AVP^n=\AVProj(k[X_0,\ldots,X_n])$, then $\pi_V$ is defined by the
homomorphism of rings
\[
X_i\longmapsto X_i^q:
k[X_0,\ldots,X_n]\longrightarrow k[X_0,\ldots,X_n]
\]
and induces the map on points
\[
(x_0:\cdots:x_n)\longmapsto(x_0^q:\cdots:x_n^q):
\AVP^n(\AVF)\longrightarrow\AVP^n(\AVF).
\]
For any regular map $\varphi:W\to V$ of varieties over $\AVF_q$, it is
obvious that $\varphi\circ\pi_W=\pi_V\circ\varphi$. Therefore, if
$V\hookrightarrow\AVP^n$ is a projective embedding of $V$, then $\pi_V$
induces the map
\[
(x_0:\ldots:x_n)\longmapsto(x_0^q:\ldots:x_n^q)
\]
on $V(\AVF)$. Thus $V(\AVF_q)$ is the set of fixed points of
$\pi_V:V(\AVF)\to V(\AVF)$.

Let $A$ be an abelian variety over $\AVF_q$. Then $\pi_A$ maps $0$ to $0$
(because $0\in V(\AVF_q)$), and so it is an endomorphism of $A$.
Recall that its characteristic polynomial $P_\pi$ is a monic polynomial
of degree $2g$, $g=\dim(A)$, with coefficients in $\AVZ$.
\AVeditorial{In the error-bound proof on parent physical page 82, the source prints $(2^g-2g-1)$ for the remaining terms; the theorem statement on physical page 81 uses $2^{2g}$. Both source expressions are retained in their own locations. This note records their printed difference and does not assert an author-issued correction.}
% END CONTEXT AV:II1.1

% BEGIN PRIMARY STATEMENT AV:II1.1
\AVsrc{2008 notes, II Theorem 1.1, physical page 81, printed page 75}
\noindent\textbf{Theorem 1.1.}
\emph{Write $P_\pi(X)=\prod_i(X-a_i)$, and let
$N_m=\lvert A(\AVF_{q^m})\rvert$. Then}
\begin{enumerate}
\item[(a)] $\displaystyle N_m=\prod_{i=1}^{2g}(1-a_i^m)$ for all
  $m\geq1$, and
\item[(b)] (Riemann hypothesis) $\lvert a_i\rvert=q^{1/2}$.
\end{enumerate}
\emph{Hence}
\[
\lvert N_m-q^{mg}\rvert
\leq 2g\cdot q^{m(g-\frac12)}
+(2^{2g}-2g-1)q^{m(g-1)}.
\]
% END PRIMARY STATEMENT AV:II1.1

% BEGIN PRIMARY PROOF AV:II1.1
\AVsrc{2008 notes, physical pages 81--82, printed pages 75--76;
complete initial proof of II1.1}
\begin{proof}
We first deduce the inequality from the preceding statements. Take $m=1$
in (a) and expand out to get
\[
\lvert A(\AVF_q)\rvert
=a_1\cdots a_{2g}
-\sum_{i=1}^{2g}a_1\cdots a_{i-1}a_{i+1}\cdots a_{2g}+\cdots.
\]
The first term on the right is an integer, and in fact a positive integer
because it is $P_\pi(0)=\deg(\pi)$, and (b) shows that it has absolute
value $q^g$. Hence it equals $q^g$ (actually, it easy to prove directly
that $\deg(\pi)=q^g$). The Riemann hypothesis shows that each term
$a_1\cdots a_{i-1}a_{i+1}\cdots a_{2g}$ has absolute value
$=q^{g-\frac12}$, and so the sum has absolute value
$\leq2g\cdot q^{g-\frac12}$. There are $(2^g-2g-1)$ terms remaining,
and each has absolute value $\leq q^{g-1}$, whence the inequality.
We first prove (a) in the case $m=1$. The kernel of
\[
\pi-\AVid:A(\AVF)\longrightarrow A(\AVF)
\]
is $A(\AVF_q)$. I claim that the map
\[
(d\pi)_0:\AVTgt_0(A)\longrightarrow\AVTgt_0(A)
\]
is zero---in fact, that this is true for any variety. In proving it, we
can replace $A$ with an open affine neighbourhood $U$, and embed $U$
into $\AVA^m$ some $m$ in such a way that $0$ maps to the origin $0$.
The map $(d\pi)_0$ on $\AVTgt_0(U)$ is the restriction of the map
$(d\pi)_0$ on $\AVTgt_0(\AVA^m)$. But
$\pi:\AVA^n\to\AVA^n$ is given by the equations $Y_i=X_i^q$, and
$d(X_i^q)=qX_i^{q-1}=0$ (in characteristic $p$). We now find that
\[
d(\pi-\AVid)_0=(d\pi)_0-(d(\AVid))_0=-1.
\]
Hence $\pi-1$ is \emph{\'etale} at the origin, and so, by homogeneity,
it is \emph{\'etale} at every point---each point in the kernel occurs
with multiplicity 1. Therefore,
\[
\lvert A(\AVF_q)\rvert=\deg(\pi-\AVid).
\]
But, from the definition of $P_\pi$, we know that
\[
\deg(\pi-\AVid)=P_\pi(1),
\]
and this is $\prod(1-a_i)$. When we replace $\pi$ with $\pi^m$ in the
above argument, we find that
\[
\lvert A(\AVF_{q^m})\rvert=P_{\pi^m}(1).
\]
Recall (10.20) that $a_1,\ldots,a_{2g}$ can be interpreted as the
eigenvalues of $\pi$ acting on $T_\ell A$. Clearly $\pi^m$ has
eigenvalues $a_1^m,\ldots,a_{2g}^m$, and so
\[
P_{\pi^m}(X)=\prod(X-a_i^m),\qquad
P_{\pi^m}(1)=\prod(1-a_i^m)
\]
which proves (a) for a general $m$.
\end{proof}

Part (b) follows from the next two lemmas.

% BEGIN SUPPORT WITHIN PRIMARY PROOF AV:II1.2
\AVsrc{2008 notes, II Lemma 1.2 and complete proof, physical pages 82--83,
printed pages 76--77}
\noindent\textbf{Lemma 1.2.}
\emph{Let $\dagger$ be the Rosati involution on $\AVEnd(A)\otimes\AVQ$
defined by a polarization of $A$; then $\pi_A^{\dagger}\circ\pi_A=q_A$.}

\begin{proof}
Let $D$ be the ample divisor on $A$ defining the polarization; thus
$\lambda(a)=[D_a-D]$. We have to show that
\[
\pi^{\vee}\circ\lambda\circ\pi=q\lambda.
\]
Recall that, on points, $\pi^{\vee}$ is the map
\[
[D']\longmapsto[\pi^*D']:\AVPic^0(A)\longrightarrow\AVPic^0(A).
\]
Let $D'$ be a divisor on $A$ (or, in fact any variety defined over
$\AVF_q$). If $D'=\AVdiv(f)$ near $\pi(P)$, then, by definition,
$\pi^*D'=\AVdiv(f\circ\pi)$ near $P$. But $\pi(P)=P$ and
$f\circ\pi=f^q$ (this was the definition of $\pi$), and
$\AVdiv(f^q)=q\AVdiv(f)$; thus $\pi^*D'=qD$. Next observe that, for
any homomorphism $\alpha:A\to A$ and any point $a$ on $A$,
\[
\alpha\circ t_a(x)=\alpha(a+x)=\alpha(a)+\alpha(x)
=t_{\alpha(a)}\circ\alpha(x).
\]
We can now prove the lemma. For any $a\in A(\AVF)$, we have
\begin{align*}
(\pi^{\vee}\circ\lambda\circ\pi)(a)
&=\pi^*[D_{\pi(a)}-D]\\
&=[\pi^*t_{\pi(a)}^*D-\pi^*D]\\
&=[(t_{\pi(a)}\circ\pi)^*D-\pi^*D]\\
&=[(\pi\circ t_a)^*D-\pi^*D]\\
&=[t_a^*\pi^*D-\pi^*D]\\
&=[t_a^*qD-qD]\\
&=q\lambda(a);
\end{align*}
as required.
\end{proof}
% END SUPPORT WITHIN PRIMARY PROOF AV:II1.2

% BEGIN SUPPORT WITHIN PRIMARY PROOF AV:II1.3
\AVsrc{2008 notes, II Lemma 1.3 and complete proof, physical pages 83--84,
printed pages 77--78, including original footnote 1}
\noindent\textbf{Lemma 1.3.}
\emph{Let $A$ be an abelian variety over a field $k$ (not necessarily
finite). Let $\alpha$ be an element of $\AVEnd(A)\otimes\AVQ$ such that
$\alpha^{\dagger}\circ\alpha$ is an integer $r$; for any root $a$ of
$P_\alpha$ in $\AVC$, $\lvert a\rvert^2=r$.}

\begin{proof}
Note that $\AVQ[\alpha]$ is a commutative ring of finite-dimension over
$\AVQ$; it is therefore an Artin ring. According to (Atiyah and
Macdonald 1969, Chapter 8)\footnote[1]{In fact, it is easy to prove this
directly. Let $f(X)$ be a monic polynomial generating the kernel of
$\AVQ[X]\to\AVQ[\alpha]$. Then
$\AVQ[X]/(f(X))\simeq\AVQ[\alpha]$, and the maximal ideals of
$\AVQ[\alpha]$ correspond to the distinct irreducible factors of $f(X)$.},
it has only finitely many prime ideals $\mathfrak m_1,\ldots,\mathfrak m_n$
each of which is also maximal, every element of $\bigcap\mathfrak m_i$
is nilpotent, and $\AVQ[\alpha]/\bigcap\mathfrak m_i$ is a product of fields
\[
\AVQ[\alpha]/\bigcap\mathfrak m_i=K_1\times\cdots\times K_n,
\qquad K_i=\AVQ[\alpha]/\mathfrak m_i.
\]
We first show that $\bigcap\mathfrak m_i=0$, i.e., that $\AVQ[\alpha]$
has no nonzero nilpotents. Note that $\AVQ[\alpha]$ is stable under the
action of $\dagger$. Let $a\ne0\in\AVQ[\alpha]$. Then
$b\overset{\mathrm{df}}=a^{\dagger}\cdot a\ne0$, because
$\AVTr(a^{\dagger}\cdot a)>0$. As $b^{\dagger}=b$,
$\AVTr(b^2)=\AVTr(b^{\dagger}\cdot b)>0$, and so $b^2\ne0$.
Similarly, $b^4\ne0$, and so on, which implies $b$ is not nilpotent,
and so neither is $a$. Any automorphism $\tau$ of $\AVQ[\alpha]$
permutes the maximal ideals $\mathfrak m_i$; it therefore permutes the
factors $K_i$, i.e., there is a permutation $\sigma$ of
$\{1,2,\ldots,n\}$ and isomorphisms $\tau_i:K_i\to K_{\sigma(i)}$
such that $\tau(a_1,\ldots,a_n)=(b_1,\ldots,b_n)$ with
$b_{\sigma(i)}=\tau_i(a_i)$. In the case that $\tau=\dagger$,
$\sigma$ must be the trivial permutation, for otherwise
$\AVTr(a^{\dagger}\cdot a)$ would not always be positive (consider
$(a_1,0,0,\ldots)$ if $\sigma(1)\ne1$). Hence $\dagger$ preserves
the factors of $\AVQ[\alpha]$, and is a positive-definite involution
on each of them. The involution $\dagger$ extends by linearity
(equivalently by continuity) to a positive-definite involution of
$\AVQ[\alpha]\otimes\AVR$. The above remarks also apply to
$\AVQ[\alpha]\otimes\AVR$: it is a product of fields, and $\dagger$
preserves each factor and is a positive-definite involution on each of
them. But now each factor is isomorphic to $\AVR$ or to $\AVC$.
The field $\AVR$ has no nontrivial automorphisms at all, and so
$\dagger$ must act on a real factor of $\AVQ[\alpha]\otimes\AVR$ as
the identity map. The field $\AVC$ has only two automorphisms of finite
order: the identity map and complex conjugation. The identity map is not
positive-definite, and so $\dagger$ must act on a complex factor as
complex conjugation. We have shown: for any homomorphism
$\rho:\AVQ[\alpha]\to\AVC$,
$\rho(\alpha^{\dagger})=\overline{\rho(\alpha)}$.
Thus, for any such homomorphism,
$r=\rho(\alpha^{\dagger}\cdot\alpha)=\lvert\rho(\alpha)\rvert^2$,
and so every root of the minimum polynomial of $\alpha$ in
$\AVQ[\alpha]/\AVQ$ has absolute value $r^{1/2}$.
Now (10.24) completes the proof.
\end{proof}
% END SUPPORT WITHIN PRIMARY PROOF AV:II1.3
% END PRIMARY PROOF AV:II1.1

% BEGIN SUPPORT CLOSING REMARK AV:II1.4
\AVsrc{2008 notes, II Remark 1.4, physical page 84, printed page 78}
\noindent\textbf{Remark 1.4.}
We have actually proved the following: $\AVQ[\pi]$ is a product of
fields, stable under the involution $\dagger$; under every map
$\tau:\AVQ[\pi]\to\AVC$, $\tau(\pi^{\dagger})=\overline{\tau(\pi)}$,
and $\lvert\tau\pi\rvert=q^{1/2}$.
% END SUPPORT CLOSING REMARK AV:II1.4

\AVeditorial{The parent text's use of $V(\AVF_q)$ in the introduction,
its $\AVA^m$/$\AVA^n$ notation and differential display in the initial
proof, and the printed conclusion $\pi^*D'=qD$ in II1.2 are retained
literally. The first proof gives the error-bound calculation for $m=1$
as printed, while the statement retains the bound for every $m$.
No additional general-$m$ calculation or additional condition on $r$ in
II1.3 has been supplied. These are transcription notes, not mathematical
repairs or a correctness assessment.}

% BEGIN SUPPORT SEPARATE VERSION AVs2022:17.3-17.4
\section*{Separate supporting source: the corrected Storrs article}
\noindent\textbf{Version boundary.}
The following is from the fixed January 2, 2022 corrected author file
$\mathrm{AVs}$, not from the 2008 course notes. Physical pages and printed
pages both read 35--36. The source retains its own numbering, notation
$l$ for the auxiliary prime, $\varphi_{\mathcal L}$, and the references
to its own (15.4). The argument is an uncounted support for the particular
trace/intersection formula omitted in the parent I14.3.

\subsection*{17. The Rosati Involution: necessary opening context}
\AVsrc{corrected Storrs article, physical/printed page 35; opening paragraphs of Section 17}
Fix a polarization $\lambda$ on $A$. As $\lambda$ is an isogeny
$A\to A^{\vee}$, it has an inverse in
$\AVHom^0(A^{\vee},A)\overset{\mathrm{def}}=
\AVHom(A^{\vee},A)\otimes\AVQ$.
The \emph{Rosati involution} on $\AVEnd^0(A)$ corresponding to
$\lambda$ is
\[
\alpha\longmapsto\alpha^{\dagger}
=\lambda^{-1}\circ\alpha^{\vee}\circ\lambda.
\]
This has the following obvious properties:
\[
(\alpha+\beta)^{\dagger}=\alpha^{\dagger}+\beta^{\dagger},\qquad
(\alpha\beta)^{\dagger}=\beta^{\dagger}\alpha^{\dagger},\qquad
a^{\dagger}=a\text{ for }a\in\AVQ.
\]
For all $a,a'\in T_l A\otimes\AVQ$, $l\ne\operatorname{char}(k)$,
\[
e_l^{\lambda}(\alpha a,a')
=e_l(\alpha a,\lambda a')
=e_l(a,\alpha^{\vee}\circ\lambda a')
=e_l^{\lambda}(a,\alpha^{\dagger}a');
\]
from which it follows that $\alpha^{\dagger\dagger}=\alpha$.

% BEGIN SUPPORT STATEMENT AVs2022:17.3
\AVsrc{corrected Storrs article, Theorem 17.3, physical/printed page 35}
\noindent\textbf{Theorem 17.3.}
\emph{The bilinear form}
\[
(\alpha,\beta)\longmapsto\AVTr(\alpha\circ\beta^{\dagger}):
\AVEnd^0(A)\times\AVEnd^0(A)\longrightarrow\AVQ
\]
\emph{is positive definite. More precisely, if
$\lambda=\varphi_{\mathcal L(D)}$, then}
\[
\AVTr(\alpha\circ\alpha^{\dagger})
=\frac{2g}{(D^g)}(D^{g-1}\cdot\alpha^*(D)).
\]
% END SUPPORT STATEMENT AVs2022:17.3

\AVsrc{corrected Storrs article, physical/printed page 36; paragraph continuing Theorem 17.3}
As $D$ is ample and $\alpha^*(D)$ is effective, the intersection number
$(D^{g-1}\cdot\alpha^*(D))$ is positive. Thus the second statement implies
the first. Clearly it suffices to prove it with $k$ algebraically closed.
We first need a lemma.

% BEGIN SUPPORT LEMMA AND PROOF AVs2022:17.4
\noindent\textbf{Lemma 17.4.}
\emph{Let $A$ be an abelian variety over an algebraically closed field,
and let $\AVZ_l(g)=\AVZ_l(1)^{\otimes g}$. Then there is a canonical
generator $\varepsilon$ of
$\AVHom(\bigwedge^{2g}(T_l A),\AVZ_l(g))$ with the following property:
if $D_1,\ldots,D_g$ are divisors on $A$ and
$e_i=e_l^{\mathcal L(D_i)}\in
\AVHom(\bigwedge^2T_l A,\AVZ_l(1))$, then
$e_1\wedge\cdots\wedge e_g$ is the multiple
$(D_1,D_2,\ldots,D_g)\varepsilon$ of $\varepsilon$.}

\begin{proof}
See Mumford 1970, \S20, Theorem 3, p. 190. (From the point of view of
\'etale cohomology, $\varepsilon$ corresponds to the canonical generator
of $H^{2g}(A,\AVZ_l(g))$, which is equal to the cohomology class of any
point on $A$. If $c_i$ is the class of $D_i$ in $H^2(A,\AVZ_l(1))$,
then the compatibility of intersection products with cup products shows
that $(D_1,\ldots,D_g)\varepsilon=c_1\cup\cdots\cup c_g$.
Consequently, the lemma follows from (15.4).
\end{proof}
% END SUPPORT LEMMA AND PROOF AVs2022:17.4

% BEGIN SUPPORT CALCULATION AVs2022:17.3
\AVsrc{corrected Storrs article, physical/printed page 36; complete calculation ending ``which completes the proof''}
\begin{proof}[Proof (of 17.3)]
From the lemma, we find that
\begin{align*}
e_l^{\lambda}\wedge\cdots\wedge e_l^{\lambda}
&=(D^g)\varepsilon,\\
e_l^{\lambda}\wedge\cdots\wedge e_l^{\lambda}
\wedge e_l^{\alpha^*(\lambda)}
&=(D^{g-1}\cdot\alpha^*(D))\varepsilon.
\end{align*}
It suffices therefore to show that, for some basis $a_1,\ldots,a_{2g}$
of $T_l A\otimes\AVQ$,
\[
\frac{\langle a_1\wedge\cdots\wedge a_{2g}\mid
e_l^{\lambda}\wedge\cdots\wedge e_l^{\lambda}
\wedge e_l^{\alpha^*(\lambda)}\rangle}
{\langle a_1\wedge\cdots\wedge a_{2g}\mid
e_l^{\lambda}\wedge\cdots\wedge e_l^{\lambda}\rangle}
=\frac{1}{2g}\AVTr(\alpha\circ\alpha^{\dagger})
\]
(see 15.4). Choose the basis $a_1,a_2,\ldots,a_{2g}$ so that
\begin{align*}
e_l^{\lambda}(a_{2i-1},a_{2i})
&=1=-e_l^{\lambda}(a_{2i},a_{2i-1}),
&&i=1,2,\ldots,g,\\
e_l^{\lambda}(a_i,a_j)&=0,&&\text{otherwise}.
\end{align*}
Let $f_1,\ldots,f_{2g}$ be the dual basis; then for $j\ne j'$,
\[
\langle a_i\wedge a_{i'}\mid f_j\wedge f_{j'}\rangle
=\begin{vmatrix}
f_j(a_i)&f_{j'}(a_i)\\
f_j(a_{i'})&f_{j'}(a_{i'})
\end{vmatrix}
=\begin{cases}
1&\text{if }i=j,\ i'=j',\\
-1&\text{if }i=j',\ i'=j,\\
0&\text{otherwise}.
\end{cases}
\]
Therefore $e_l^{\lambda}=\sum_{i=1}^g f_{2i-1}\wedge f_{2i}$, and so
$e_l^{\lambda}\wedge\cdots\wedge e_l^{\lambda}
=g!(f_1\wedge\cdots\wedge f_{2g})$. Thus
\[
\langle a_1\wedge\cdots\wedge a_{2g}\mid
e_l^{\lambda}\wedge\cdots\wedge e_l^{\lambda}\rangle
=\langle a_1\wedge\cdots\wedge a_{2g}\mid
g!(f_1\wedge\cdots\wedge f_{2g})\rangle=g!.
\]
Similarly,
\begin{align*}
&\langle a_1\wedge\cdots\wedge a_{2g}\mid
e_l^{\lambda}\wedge\cdots\wedge e_l^{\lambda}
\wedge e_l^{\alpha^*(\mathcal L)}\rangle\\
&\qquad=(g-1)!\sum_{i=1}^g
e_l^{\lambda}(\alpha a_{2i-1},\alpha a_{2i})\\
&\qquad=\frac{(g-1)!}{2}\sum
\bigl(e_l^{\lambda}(a_{2i-1},\alpha^{\dagger}\alpha a_{2i})
+e_l^{\lambda}(\alpha^{\dagger}\alpha a_{2i-1},a_{2i})\bigr)\\
&\qquad=\frac{g!}{2g}\AVTr(\alpha^{\dagger}\alpha),
\end{align*}
which completes the proof.
\end{proof}
% END SUPPORT CALCULATION AVs2022:17.3

\AVeditorial{The source's explicit algebraic-closure reduction and the
auxiliary-prime condition $l\ne\operatorname{char}(k)$ are retained;
there is no characteristic-zero or principal-polarization restriction
added here. Lemma 17.4 retains both its Mumford reference and its printed
cup-product explanation with the same article's (15.4). Those background
interfaces are not re-proved. The source's unclosed explanatory parenthesis
in that proof, its use of $e_l^{\alpha^*(\lambda)}$ followed later by
$e_l^{\alpha^*(\mathcal L)}$, and the final unindexed $\sum$ are kept
as printed, without mathematical emendation.}
% END SUPPORT SEPARATE VERSION AVs2022:17.3-17.4

\end{document}
